% =====================================================================
% pitch.tex - Pitch Beamer (10 min) IoT Edge Bridge
% Projet TS6 - Noel Junior Yando Fotso - EPHEC 2025-2026
% Compilation : latexmk -xelatex -shell-escape pitch.tex
% =====================================================================
% --- Toggle mode notes (presentation dual-ecran) --------------------
% Par defaut : 0 = PDF de projection, slides seules.
% Pour generer le PDF avec notes (pour pdfpc / Skim / Acrobat dual-ecran) :
%   make pitch-notes
% qui passe -pretex='\def\PRESENTERNOTES{1}' a latexmk.
% \def car \providecommand n'existe pas encore avant \documentclass.
\ifx\PRESENTERNOTES\undefined\def\PRESENTERNOTES{0}\fi

\documentclass[10pt,aspectratio=169,xcolor={table}]{beamer}

% --- Theme -----------------------------------------------------------
% Theme Metropolis si disponible, sinon fallback "default" (commente).
\IfFileExists{beamerthememetropolis.sty}{%
  \usetheme{metropolis}%
}{%
  % Fallback : theme par defaut. Decommentez pour forcer un autre theme.
  \usetheme{default}%
  % \usetheme{Madrid}
}

% --- Encodage / langue / fontes --------------------------------------
\usepackage[utf8]{inputenc}
\usepackage[T1]{fontenc}
\usepackage[french]{babel}
\usepackage{csquotes}

\usepackage{ifxetex,ifluatex}
\newif\ifxetexorluatex
\ifxetex\xetexorluatextrue\else
  \ifluatex\xetexorluatextrue\else\xetexorluatexfalse\fi
\fi
\ifxetexorluatex
  \usepackage{fontspec}
  \IfFontExistsTF{Fira Sans}{\setsansfont{Fira Sans}}{}
  \IfFontExistsTF{Fira Mono}{\setmonofont{Fira Mono}}{}
\fi

% --- Graphiques / couleurs / TikZ ------------------------------------
\usepackage{graphicx}
\usepackage{xcolor}
\definecolor{bridgeblue}{HTML}{1F3A68}
\definecolor{bridgeteal}{HTML}{0FA4AF}
\definecolor{bridgegray}{HTML}{4A4A4A}
\definecolor{bridgealert}{HTML}{C62828}
\definecolor{bridgesoft}{HTML}{F2F6FB}
\definecolor{bridgeinfo}{HTML}{E3F4F4}
\definecolor{bridgegreen}{HTML}{2E7D32}
\definecolor{bridgegreensoft}{HTML}{E8F5E9}

\usepackage{tikz}
\usetikzlibrary{shapes.geometric, arrows.meta, positioning, fit, backgrounds, calc}

% Styles TikZ partages (necessaires pour les figures includes du projet)
\tikzset{
  bridgeblock/.style={%
    rectangle, rounded corners=2pt,
    draw=bridgeblue, thick,
    fill=bridgeblue!8,
    text=bridgeblue,
    align=center,
    minimum width=28mm, minimum height=10mm,
    inner sep=4pt,
    font=\small\sffamily
  },
  bridgeflow/.style={%
    -{Stealth[length=3mm,width=2mm]},
    draw=bridgeteal, thick
  }
}

\usepackage{booktabs}
\usepackage{tabularx}
\usepackage{listings}
\usepackage{enumitem}
\usepackage{colortbl}
\usepackage{array}

% --- Icones FontAwesome 5 (avec fallback silencieux si absent) -------
\IfFileExists{fontawesome5.sty}{%
  \usepackage{fontawesome5}%
}{%
  % Fallback : on definit des macros vides pour ne pas casser la compilation.
  \newcommand{\faHospital}{}\newcommand{\faHeartbeat}{}\newcommand{\faTint}{}%
  \newcommand{\faLungs}{}\newcommand{\faVial}{}\newcommand{\faCogs}{}%
  \newcommand{\faUserShield}{}\newcommand{\faLock}{}\newcommand{\faServer}{}%
  \newcommand{\faDatabase}{}\newcommand{\faNetworkWired}{}\newcommand{\faPlug}{}%
  \newcommand{\faExclamationTriangle}{}\newcommand{\faCheckCircle}{}%
  \newcommand{\faClock}{}\newcommand{\faWifi}{}\newcommand{\faStethoscope}{}%
  \newcommand{\faRoute}{}\newcommand{\faProjectDiagram}{}\newcommand{\faPython}{}%
  \newcommand{\faDocker}{}\newcommand{\faFile}{}\newcommand{\faRandom}{}%
  \newcommand{\faChartLine}{}\newcommand{\faBalanceScale}{}%
  \newcommand{\faQuestionCircle}{}\newcommand{\faUserMd}{}\newcommand{\faGavel}{}%
  \newcommand{\faKey}{}\newcommand{\faRedo}{}\newcommand{\faClipboardCheck}{}%
}

% --- Macro zebra : alternance des lignes (booktabs friendly) ---------
% \zebra avant un tableau : grise les lignes paires depuis la 2e ligne.
\newcommand{\zebra}{\rowcolors{2}{bridgesoft}{white}}

% --- Styles listings -------------------------------------------------
\lstdefinestyle{bridgebase}{%
  basicstyle=\ttfamily\scriptsize,
  breaklines=true,
  columns=fullflexible,
  keepspaces=true,
  showstringspaces=false,
  frame=single,
  framerule=0.4pt,
  rulecolor=\color{bridgegray!40},
  backgroundcolor=\color{bridgegray!5},
}
\lstdefinestyle{bridgejson}{%
  style=bridgebase,
  basicstyle=\ttfamily\scriptsize\color{bridgeblue!90!black},
  keywordstyle=\color{bridgeteal}\bfseries,
  stringstyle=\color{bridgealert!80!black},
  morekeywords={resourceType,code,coding,system,valueQuantity,value,unit,subject,effectiveDateTime,status},
  morestring=[b]",
}

% --- Couleurs Beamer projet -----------------------------------------
\setbeamercolor{frametitle}{bg=bridgeblue,fg=white}
\setbeamercolor{title}{fg=bridgeblue}
\setbeamercolor{structure}{fg=bridgeblue}
\setbeamercolor{alerted text}{fg=bridgealert}
\setbeamercolor{block title}{bg=bridgeteal,fg=white}
\setbeamercolor{block body}{bg=bridgeteal!8,fg=black}

% --- Desactivation explicite des slides de section automatiques ------
\metroset{sectionpage=none}

% --- Encadres colores (warn / info / key / note) ---------------------
\newenvironment{bridgewarn}[1][Attention]{%
  \begin{beamercolorbox}[wd=\textwidth,sep=2mm,rounded=true]{block body}%
    \footnotesize\textbf{\color{bridgealert}#1.}\par\vspace{1mm}%
}{%
  \end{beamercolorbox}%
}
\newenvironment{bridgeinfo}[1][Note]{%
  \begin{beamercolorbox}[wd=\textwidth,sep=2mm,rounded=true]{block body}%
    \footnotesize\textbf{\color{bridgeblue}#1.}\par\vspace{1mm}%
}{%
  \end{beamercolorbox}%
}
\newenvironment{bridgekey}[1][Cle]{%
  \begin{beamercolorbox}[wd=\textwidth,sep=2mm,rounded=true]{block body}%
    \footnotesize\textbf{\color{bridgeteal}#1.}\par\vspace{1mm}%
}{%
  \end{beamercolorbox}%
}

% --- Footer custom 4 zones : auteur | partie | chapitre | nav + n slide
% Variables sans @ pour eviter les problemes de catcode hors preambule.
\newcommand{\PartCurrent}{}
\newcommand{\ChCurrent}{}
\newcommand{\partref}[1]{\renewcommand{\PartCurrent}{#1}}
\newcommand{\chref}[1]{\renewcommand{\ChCurrent}{#1}}

% --- Boutons de navigation cliquables --------------------------------
% prev/next utilisent \hyperlink* Beamer (marchent dans TOUT lecteur PDF
% qui supporte les liens). Le bouton plein ecran utilise \Acrobatmenu
% (hyperref) qui marche en Adobe Acrobat (Windows/Mac) et la plupart
% des lecteurs PDF natifs (pas dans les browsers).
\newcommand{\NavBtn}[1]{\textcolor{bridgeblue}{\large#1}}

\setbeamertemplate{footline}{%
  \leavevmode%
  \hbox{%
    \begin{beamercolorbox}[wd=0.24\paperwidth,ht=2.8ex,dp=1.2ex,leftskip=2ex]{author in head/foot}%
      \usebeamerfont{author in head/foot}\insertshortauthor%
    \end{beamercolorbox}%
    \begin{beamercolorbox}[wd=0.30\paperwidth,ht=2.8ex,dp=1.2ex,center]{title in head/foot}%
      \usebeamerfont{title in head/foot}\textcolor{bridgegray}{\PartCurrent}%
    \end{beamercolorbox}%
    \begin{beamercolorbox}[wd=0.18\paperwidth,ht=2.8ex,dp=1.2ex,center]{title in head/foot}%
      \usebeamerfont{title in head/foot}\textcolor{bridgegray}{\ChCurrent}%
    \end{beamercolorbox}%
    \begin{beamercolorbox}[wd=0.28\paperwidth,ht=2.8ex,dp=1.2ex,rightskip=2ex,right]{title in head/foot}%
      \usebeamerfont{title in head/foot}%
      \hyperlinkslideprev{\NavBtn{\faChevronCircleLeft}}\,%
      \Acrobatmenu{FullScreen}{\NavBtn{\faExpand}}\,%
      \hyperlinkslidenext{\NavBtn{\faChevronCircleRight}}\quad%
      \textcolor{bridgegray}{\insertframenumber/\inserttotalframenumber}%
    \end{beamercolorbox}%
  }%
}

\setbeamertemplate{navigation symbols}{}

% --- Mode notes a droite (pour pdfpc/Skim/Acrobat dual-ecran) -------
\ifnum\PRESENTERNOTES=1
  \usepackage{pgfpages}
  \setbeameroption{show notes on second screen=right}
\fi

% --- Hyperref : ouvrir en plein ecran + bookmarks --------------------
% beamer charge deja hyperref, on customise juste les options.
\hypersetup{%
  pdftitle={IoT Edge Bridge - Pitch TS6 EPHEC},
  pdfauthor={Noel Junior Yando Fotso},
  pdfsubject={Passerelle d'interoperabilite Edge pour le BIHR},
  pdfkeywords={FHIR, HL7, IoT, BIHR, Edge, Interoperabilite},
  pdfpagemode=FullScreen,
  pdfstartview=Fit,
  bookmarks=true,
  bookmarksopen=true
}

% --- Metadonnees -----------------------------------------------------
\title{IoT Edge Bridge}
\subtitle{Du chaos des dispositifs au continuum BIHR}
\author[Yando Fotso]{No\"el Junior Yando Fotso}
\institute{EPHEC Brussels --- TS6}
\date{2 mars 2027}

% =====================================================================
\begin{document}

% =====================================================================
% F1. COUVERTURE
% =====================================================================
\partref{}
\chref{}
\begin{frame}[plain]
  \titlepage
  \note{Bonjour a tous. Je vais vous raconter une histoire en 10 minutes. Celle d'un parc de dispositifs medicaux qui ne se parlent pas, et d'une passerelle qui les reconcilie. Demo en live a la fin.}
\end{frame}

% =====================================================================
% F2. PREAMBULE NARRATIF
% =====================================================================
\partref{Partie 1/5 : Le contexte}
\chref{Ch.1}
\begin{frame}{Le contexte}
\begin{itemize}[leftmargin=*, itemsep=0.4em]
  \item \textcolor{bridgealert}{\faHeartbeat}~\textbf{Des milliers} de mesures vitales par jour : moniteurs, pompes, glucom\`etres, ventilateurs.
  \item \textcolor{bridgeblue}{\faHospital}~\textbf{Plan eSant\'e 2025-2027} : toutes ces mesures attendues dans le BIHR.
  \item \textcolor{bridgealert}{\faExclamationTriangle}~\textbf{Verrou} : l'interop\'erabilit\'e, pas la volont\'e.
  \item \textcolor{bridgealert}{\faGavel}~\textbf{Contrainte} : toucher au mat\'eriel certifi\'e casse le marquage CE.
\end{itemize}

\vspace{1em}
\begin{bridgekey}[\faKey~Le mandat]
Faire arriver les mesures au BIHR \textbf{sans toucher aux dispositifs}.
\end{bridgekey}
\note{Slide 2. On commence par poser le contexte humain et professionnel. Ce n'est pas un probleme theorique : ces mesures existent, elles ne se rendent pas au DPI, et c'est mauvais pour le patient.}
\end{frame}

% =====================================================================
% F3. PROCESSUS METIER
% =====================================================================
\partref{Partie 1/5 : Le contexte}
\chref{Ch.1}
\begin{frame}{Le processus m\'etier qu'on automatise}
  \begin{center}
    \resizebox{!}{0.74\textheight}{% =====================================================================
% figures/fig-processus-metier.tex
% Diagramme de processus metier BPMN simplifie : source -> bridge -> DPI -> soignant.
% Swimlane horizontale a 4 lignes.
% =====================================================================
\begin{tikzpicture}[
  font=\sffamily,
  every node/.style={font=\footnotesize\sffamily},
  >={Stealth[length=2.5mm]},
  thick,
  source/.style={
    rectangle, rounded corners=2pt,
    draw=bridgealert, fill=bridgealert!10, text=bridgealert,
    minimum width=2.4cm, minimum height=0.8cm, align=center,
    inner sep=2pt, font=\footnotesize\sffamily
  },
  bridge/.style={
    rectangle, rounded corners=2pt,
    draw=bridgeblue, fill=bridgeblue!10, text=bridgeblue,
    minimum width=2.4cm, minimum height=0.8cm, align=center,
    inner sep=2pt, font=\footnotesize\sffamily
  },
  dpi/.style={
    rectangle, rounded corners=2pt,
    draw=bridgeteal, fill=bridgeteal!12, text=bridgeteal,
    minimum width=2.4cm, minimum height=0.8cm, align=center,
    inner sep=2pt, font=\footnotesize\sffamily
  },
  user/.style={
    rectangle, rounded corners=2pt,
    draw=bridgegreen, fill=bridgegreen!12, text=bridgegreen,
    minimum width=2.4cm, minimum height=0.8cm, align=center,
    inner sep=2pt, font=\footnotesize\sffamily
  },
  decision/.style={
    diamond, draw=bridgeblue, fill=bridgeblue!8, text=bridgeblue,
    minimum width=1.4cm, minimum height=1.0cm, inner sep=1pt,
    font=\scriptsize\sffamily, aspect=1.6
  },
  flow/.style={-{Stealth[length=3mm]}, thick, draw=bridgeblue!80},
  errflow/.style={-{Stealth[length=3mm]}, thick, dashed, draw=bridgealert}
]

% --- Lignes (swim lanes) ---------------------------------------------
\foreach \y/\name in {3.6/Source, 2.4/Bridge, 1.2/DPI, 0/Soignant}
  \node[anchor=east, font=\footnotesize\bfseries\sffamily, color=bridgegray]
    at (-0.2, \y) {\name};

% Lignes de separation
\foreach \y in {3.0, 1.8, 0.6}
  \draw[bridgegray!30, dashed] (-0.1,\y) -- (16.5,\y);

% --- Bandeau Source --------------------------------------------------
\node[source] (s1) at (1.2, 3.6) {Dispositif IoT\\ou logiciel Legacy};
\node[font=\scriptsize\itshape, color=bridgegray]
  at (s1.south)[anchor=north, yshift=-1pt] {emet trame};

% --- Bandeau Bridge --------------------------------------------------
\node[bridge] (b1) at (4.0, 2.4) {Capture\\Adapter};
\node[bridge] (b2) at (6.4, 2.4) {Parsing\\Mapping};
\node[decision] (d1) at (8.8, 2.4) {Valide ?};
\node[bridge] (b3) at (11.2, 2.4) {Transport\\mTLS};

% Quarantaine en cas d'erreur
\node[bridge, fill=bridgealert!10, draw=bridgealert, text=bridgealert]
  (q) at (8.8, 3.6) {Quarantaine\\+ alerte};

% --- Bandeau DPI -----------------------------------------------------
\node[dpi] (dpi1) at (13.6, 1.2) {DPI hospitalier\\persiste FHIR};

% --- Bandeau Soignant ------------------------------------------------
\node[user] (u1) at (15.6, 0) {Consultation\\au DPI};

% --- Fleches du flux nominal -----------------------------------------
\draw[flow] (s1.south) -- (b1.north -| s1.south);
\draw[flow] (s1.south) |- (b1.west);
\draw[flow] (b1.east) -- (b2.west);
\draw[flow] (b2.east) -- (d1.west);
\draw[flow] (d1.east) -- (b3.west) node[midway, above, font=\scriptsize] {oui};
\draw[flow] (b3.south) |- (dpi1.west);
\draw[flow] (dpi1.south) |- (u1.west);

% --- Fleche d'erreur -------------------------------------------------
\draw[errflow] (d1.north) -- (q.south) node[midway, right, font=\scriptsize, color=bridgealert] {non};

\end{tikzpicture}
}\\[1mm]
    {\itshape\scriptsize 4 lanes (Source, Bridge, DPI, Soignant). 8 t\^aches, 1 d\'ecision, flux s\'equence et message en pointill\'e.}
  \end{center}
  \note{BPMN 2.0 standard. 4 lanes, evenements debut/fin, taches, gateway de validation, branche de quarantaine sur erreur, flux de message en pointille entre lanes, annotations dashed pour expliquer les contraintes. Aujourd'hui, ces 8 etapes sont 100 pourcent manuelles. Notre passerelle les automatise.}
\end{frame}

% =====================================================================
% F4. VUE SCHEMATIQUE DES ECHANGES
% =====================================================================
\partref{Partie 2/5 : Le syst\`eme}
\chref{Ch.3}
\begin{frame}{Vue sch\'ematique des \'echanges}
  \begin{center}
    \resizebox{0.92\textwidth}{!}{% =====================================================================
% figures/fig-chaos-ordre.tex
% Vision narrative : du chaos des dispositifs au continuum BIHR.
% Trois zones : chaos (rouge), bridge (bleu, 4 couches avec description),
% ordre (vert).
% =====================================================================
\begin{tikzpicture}[
  font=\sffamily,
  every node/.style={font=\small\sffamily},
  >={Stealth[length=2.5mm]},
  thick,
  device/.style={
    rectangle, rounded corners=2pt,
    draw=bridgealert, fill=bridgealert!10, text=bridgealert,
    minimum width=2.4cm, minimum height=0.7cm,
    inner sep=2pt, font=\small\sffamily
  },
  proto/.style={
    font=\scriptsize\itshape, color=bridgegray,
    inner sep=1pt
  },
  layer/.style={
    rectangle, rounded corners=2pt,
    draw=bridgeblue, fill=bridgeblue!10, text=bridgeblue,
    minimum width=2.0cm, minimum height=0.7cm,
    inner sep=2pt, font=\small\bfseries\sffamily
  },
  layerdesc/.style={
    text=bridgegray, align=left,
    inner sep=1pt,
    font=\scriptsize\sffamily
  },
  target/.style={
    rectangle, rounded corners=2pt,
    draw=bridgegreen, fill=bridgegreen!10, text=bridgegreen,
    minimum width=3.2cm, minimum height=0.7cm,
    inner sep=2pt, font=\small\sffamily
  },
  zonebox/.style={rounded corners=4pt, thick, fill opacity=1},
  flowarrow/.style={-{Stealth[length=3mm]}, very thick}
]

% =====================================================================
% Zone 1 : Chaos (gauche) - dispositifs alignes verticalement
% =====================================================================
\node[zonebox, draw=bridgealert!70, dashed, fill=bridgealert!4,
      minimum width=5.0cm, minimum height=6.4cm,
      anchor=north west, label={[anchor=north, font=\small\bfseries,
                                 color=bridgealert, yshift=-4pt]north:Chaos cote sources}]
      (chaoszone) at (0, 6.2) {};

\node[device] (d1) at (1.4, 4.8) {Moniteur};
\node[device] (d2) at (1.4, 4.05) {Glucometre};
\node[device] (d3) at (1.4, 3.30) {Pompe};
\node[device] (d4) at (1.4, 2.55) {Logiciel labo};
\node[device] (d5) at (1.4, 1.80) {DMS Legacy};
\node[device] (d6) at (1.4, 1.05) {Export Excel};

\node[proto, anchor=west] at ([xshift=2pt]d1.east) {RS-232};
\node[proto, anchor=west] at ([xshift=2pt]d2.east) {BLE/JSON};
\node[proto, anchor=west] at ([xshift=2pt]d3.east) {MEDIBUS};
\node[proto, anchor=west] at ([xshift=2pt]d4.east) {HL7v2};
\node[proto, anchor=west] at ([xshift=2pt]d5.east) {XML maison};
\node[proto, anchor=west] at ([xshift=2pt]d6.east) {CSV/XLSX};

% =====================================================================
% Zone 2 : Bridge (centre) - 4 couches AVEC description courte a droite
% =====================================================================
\node[zonebox, draw=bridgeblue, fill=bridgeblue!4,
      minimum width=5.4cm, minimum height=6.4cm,
      anchor=north west, label={[anchor=north, font=\small\bfseries,
                                 color=bridgeblue, yshift=-4pt]north:IoT Edge Bridge}]
      (bridgezone) at (6.0, 6.2) {};

% Coordonnees X : couche a x=7.0, description a x=8.6
\def\xL{7.0}
\def\xD{8.6}

\node[layer, anchor=center] (adapter) at (\xL, 4.95) {Adapter};
\node[layerdesc, anchor=west] at (\xD, 4.95)
  {Capte le canal :\\port serie, MQTT,\\TCP, fichier};

\node[layer, anchor=center] (parser) at (\xL, 3.85) {Parser};
\node[layerdesc, anchor=west] at (\xD, 3.85)
  {Decode la trame :\\HL7v2, JSON,\\MEDIBUS, CSV};

\node[layer, anchor=center] (mapper) at (\xL, 2.75) {Mapper};
\node[layerdesc, anchor=west] at (\xD, 2.75)
  {Codifie en LOINC,\\ecrit le Bundle\\FHIR R4};

\node[layer, anchor=center] (transp) at (\xL, 1.65) {Transport};
\node[layerdesc, anchor=west] at (\xD, 1.65)
  {Publie en mTLS\\avec retry et\\fallback SQLite};

\draw[bridgeflow] (adapter.south) -- (parser.north);
\draw[bridgeflow] (parser.south)  -- (mapper.north);
\draw[bridgeflow] (mapper.south)  -- (transp.north);

\node[font=\scriptsize\itshape, color=bridgegray, anchor=north]
  at (8.2, 1.05) {Profils JSON + plugins};

% =====================================================================
% Zone 3 : Ordre (droite) - destinataires alignes verticalement
% =====================================================================
\node[zonebox, draw=bridgegreen!70, fill=bridgegreen!4,
      minimum width=4.6cm, minimum height=6.4cm,
      anchor=north west, label={[anchor=north, font=\small\bfseries,
                                 color=bridgegreen, yshift=-4pt]north:Ordre cote plateforme}]
      (ordrezone) at (12.0, 6.2) {};

\node[target] (fhir) at (14.3, 4.8) {HAPI FHIR R4};
\node[target] (dpi)  at (14.3, 3.9) {DPI hospitalier};
\node[target] (bihr) at (14.3, 3.0) {MetaHub / BIHR};
\node[target] (cli)  at (14.3, 2.1) {Clinicien};

\draw[bridgeflow] (fhir.south) -- (dpi.north);
\draw[bridgeflow] (dpi.south)  -- (bihr.north);
\draw[bridgeflow] (bihr.south) -- (cli.north);

% =====================================================================
% Fleches inter-zones, parfaitement horizontales au niveau du milieu
% =====================================================================
\draw[flowarrow, draw=bridgegray]
  (chaoszone.east |- 0,3.5) -- (bridgezone.west |- 0,3.5)
  node[midway, above, font=\scriptsize\itshape, color=bridgegray] {capture};

\draw[flowarrow, draw=bridgegray]
  (bridgezone.east |- 0,3.5) -- (ordrezone.west |- 0,3.5)
  node[midway, above, font=\scriptsize\itshape, color=bridgegray] {publication};

\end{tikzpicture}
}\\[2mm]
    {\itshape\small \`A gauche, des sources h\'et\'erog\`enes. Au centre, un seul middleware. \`A droite, le DPI et le BIHR uniformis\'es.}
  \end{center}
  \note{A gauche le chaos, six dialectes differents. Au centre la passerelle, traducteur unique. A droite l'ordre, un seul format standardise vers le DPI et le BIHR.}
\end{frame}

% =====================================================================
% F5. STANDARDS DE L'INTEROPERABILITE
% =====================================================================
\partref{Partie 2/5 : Le syst\`eme}
\chref{Ch.3}
\begin{frame}{Les standards qui font le pont}
  \begin{center}
    \scriptsize
    \zebra
    \begin{tabularx}{\linewidth}{@{}c l l X@{}}
      \toprule
      & \textbf{Standard} & \textbf{Position} & \textbf{En une phrase} \\
      \midrule
      \textcolor{bridgealert}{\faFile}      & HL7 v2   & Entr\'ee Legacy & Format texte \`a barres verticales utilis\'e par 80~\% du parc hospitalier mondial. \\
      \textcolor{bridgeteal}{\faRoute}         & FHIR R4  & Sortie moderne  & Norme HL7 r\'ecente, ressources REST/JSON, lisible par les DPI 2020+. \\
      \textcolor{bridgeblue}{\faProjectDiagram}& LOINC    & Codification    & Dictionnaire universel de codes pour mesures et observations cliniques. \\
      \textcolor{bridgegreen}{\faNetworkWired} & IHE PCD  & Cadre intero.   & Profils IHE de connexion des dispositifs m\'edicaux au DPI. \\
      \textcolor{bridgeblue}{\faHospital}      & KMEHR    & Belgique        & Format XML belge des \'echanges hubs r\'egionaux et MetaHub. \\
      \textcolor{bridgegray}{\faCogs}          & DICOM    & Imagerie        & Standard mondial de l'imagerie m\'edicale (hors scope d\'emo). \\
      \bottomrule
    \end{tabularx}
  \end{center}
  \note{On ne reinvente rien. HL7 v2 pour absorber le legacy. FHIR R4 pour publier moderne. LOINC pour les codes. IHE PCD comme cadre. KMEHR et DICOM en roadmap. Chacun a sa place, chacun fait son metier.}
\end{frame}

% =====================================================================
% F6. ARCHITECTURE TECHNIQUE : FONDEMENTS (4 COUCHES)
% =====================================================================
\partref{Partie 3/5 : L'architecture}
\chref{Ch.4}
\begin{frame}{L'architecture : quatre couches enfichables}
  \begin{center}
    \resizebox{0.88\textwidth}{!}{% =====================================================================
% figures/fig-pipeline-4couches.tex
% Pipeline fonctionnel : Adapter -> Parser -> Mapper -> Transport -> HAPI
% Inputs varies a gauche, sortie FHIR a droite. Bandeau plugins au-dessus.
% =====================================================================
\begin{tikzpicture}[
  font=\sffamily,
  every node/.style={font=\small\sffamily},
  >={Stealth[length=2.5mm]},
  thick,
  layer/.style={
    rectangle, rounded corners=3pt,
    draw=bridgeblue, fill=bridgeblue!10, text=bridgeblue,
    minimum width=2.4cm, minimum height=4.2cm,
    align=center,
    font=\small\bfseries\sffamily
  },
  altlayer/.style={
    rectangle, rounded corners=3pt,
    draw=bridgeteal, fill=bridgeteal!10, text=bridgeteal,
    minimum width=2.4cm, minimum height=4.2cm,
    align=center,
    font=\small\bfseries\sffamily
  },
  outbox/.style={
    rectangle, rounded corners=3pt,
    draw=bridgegreen, fill=bridgegreen!15, text=bridgegreen,
    minimum width=2.0cm, minimum height=1.0cm,
    align=center,
    font=\small\bfseries\sffamily
  },
  inputlbl/.style={
    anchor=east, color=bridgegray,
    font=\footnotesize\itshape\sffamily,
    inner sep=2pt
  },
  caption/.style={
    align=center, font=\scriptsize\itshape\sffamily,
    color=bridgegray, inner sep=1pt
  },
  flowarrow/.style={-{Stealth[length=3mm]}, thick, draw=bridgeblue}
]

% =====================================================================
% Etiquettes des entrees (a gauche)
% =====================================================================
\node[inputlbl] (lblTcp) at (-0.2, 4.5) {TCP / HL7v2};
\node[inputlbl] (lblBle) at (-0.2, 3.7) {BLE / MQTT};
\node[inputlbl] (lblFil) at (-0.2, 2.9) {Fichier XML};
\node[inputlbl] (lblRs)  at (-0.2, 2.1) {RS-232};

% =====================================================================
% Couche 1 : Adapter (teal pour distinguer la frontiere d'entree)
% =====================================================================
\node[altlayer] (adapter) at (1.4, 3.3) {Adapter\\Layer};

% Fleches d'entree
\draw[flowarrow] (lblTcp.east) -- (adapter.west |- lblTcp.east);
\draw[flowarrow] (lblBle.east) -- (adapter.west |- lblBle.east);
\draw[flowarrow] (lblFil.east) -- (adapter.west |- lblFil.east);
\draw[flowarrow] (lblRs.east)  -- (adapter.west |- lblRs.east);

% =====================================================================
% Couche 2 : Parser
% =====================================================================
\node[layer] (parser) at (4.6, 3.3) {Parser\\Layer};
\draw[flowarrow] (adapter.east) -- (parser.west)
  node[midway, above, font=\scriptsize\itshape, color=bridgegray, yshift=2pt]
    {trame brute};

% =====================================================================
% Couche 3 : Mapper
% =====================================================================
\node[layer] (mapper) at (7.8, 3.3) {Mapper\\Layer};
\draw[flowarrow] (parser.east) -- (mapper.west)
  node[midway, above, font=\scriptsize\itshape, color=bridgegray, yshift=2pt]
    {objet};

% =====================================================================
% Couche 4 : Transport
% =====================================================================
\node[altlayer] (transp) at (11.0, 3.3) {Transport\\Layer};
\draw[flowarrow] (mapper.east) -- (transp.west)
  node[midway, above, font=\scriptsize\itshape, color=bridgegray, yshift=2pt]
    {FHIR};

% =====================================================================
% Sortie HAPI
% =====================================================================
\node[outbox] (hapi) at (13.9, 3.3) {HAPI FHIR};
\draw[flowarrow] (transp.east) -- (hapi.west)
  node[midway, above, font=\scriptsize\itshape, color=bridgegray, yshift=2pt]
    {mTLS};

% =====================================================================
% Annotations en bas (bien sous chaque couche, espace suffisant)
% =====================================================================
\node[caption] at (1.4, 0.7) {drivers\\d'entree};
\node[caption] at (4.6, 0.7) {profils JSON\\decodage};
\node[caption] at (7.8, 0.7) {LOINC\\Observation FHIR};
\node[caption] at (11.0, 0.7) {retry, fallback\\file SQLite};

% =====================================================================
% Bandeau plugins (au-dessus, couvrant les 3 couches extensibles)
% =====================================================================
\node[draw=bridgegray, dashed, fill=bridgegray!8, rounded corners=2pt,
      minimum width=8.4cm, minimum height=0.7cm,
      anchor=south, font=\scriptsize\itshape\sffamily, color=bridgegray]
  at (4.6, 6.0)
  {Plugins communautaires : extensions Adapter, Parser, Mapper};

% Petites fleches verticales du bandeau vers les 3 couches
\foreach \x in {1.4, 4.6, 7.8} {
  \draw[->, dashed, color=bridgegray, thin] (\x, 5.95) -- (\x, 5.45);
}

\end{tikzpicture}
}\\[2mm]
    {\itshape\small Adapter capte le canal. Parser d\'ecode la langue. Mapper codifie. Transport publie en s\'ecurit\'e.}
  \end{center}
  \note{Quatre couches. Adapter capte le canal. Parser decode la langue. Mapper codifie en LOINC et FHIR. Transport publie de maniere securisee. Pour ajouter un dispositif, on repond a trois questions dans un fichier JSON. Aucune ligne de code.}
\end{frame}

% =====================================================================
% F7. ARCHITECTURE TECHNIQUE : TOPOLOGIE DOCKER
% =====================================================================
\partref{Partie 3/5 : L'architecture}
\chref{Ch.4}
\begin{frame}{Topologie Docker : tout sur un poste}
  \begin{center}
    \resizebox{!}{0.72\textheight}{% =====================================================================
% figures/fig-topologie-docker.tex
% Topologie Docker Compose : 4 simulateurs + bridge + HAPI + MQTT broker.
% Layout en 4 niveaux verticaux : sims, bus mqtt+volume, bridge+dashboard, hapi
% =====================================================================
\begin{tikzpicture}[
  font=\sffamily,
  every node/.style={font=\small\sffamily},
  >={Stealth[length=2.5mm]},
  thick,
  sim/.style={
    rectangle, rounded corners=3pt,
    draw=bridgealert, fill=bridgealert!10, text=bridgealert,
    minimum width=2.6cm, minimum height=1.0cm,
    align=center, font=\small\sffamily
  },
  bus/.style={
    rectangle, rounded corners=3pt,
    draw=bridgeteal, fill=bridgeteal!12, text=bridgeteal,
    minimum width=2.6cm, minimum height=0.9cm,
    align=center, font=\small\sffamily
  },
  vol/.style={
    rectangle, rounded corners=3pt,
    draw=bridgegray, dashed, fill=bridgegray!8,
    minimum width=2.6cm, minimum height=0.9cm,
    align=center, font=\small\sffamily, text=bridgegray
  },
  bridgecore/.style={
    rectangle, rounded corners=4pt,
    draw=bridgeblue, fill=bridgeblue!12, text=bridgeblue,
    minimum width=6.6cm, minimum height=1.6cm,
    align=center, font=\bfseries\sffamily
  },
  dashbox/.style={
    rectangle, rounded corners=3pt,
    draw=bridgeteal, fill=bridgeteal!15, text=bridgeteal,
    minimum width=2.6cm, minimum height=1.6cm,
    align=center, font=\small\bfseries\sffamily
  },
  hapi/.style={
    rectangle, rounded corners=3pt,
    draw=bridgegreen, fill=bridgegreen!15, text=bridgegreen,
    minimum width=3.0cm, minimum height=1.0cm,
    align=center, font=\bfseries\sffamily
  },
  protolbl/.style={
    font=\scriptsize\itshape\sffamily, color=bridgegray, inner sep=1pt
  },
  flowarrow/.style={-{Stealth[length=3mm]}, thick, draw=bridgeblue!80}
]

% =====================================================================
% Cadre du reseau Docker
% =====================================================================
\node[draw=bridgeblue, dashed, thick, fill=bridgesoft,
      rounded corners=4pt,
      minimum width=15.6cm, minimum height=10.0cm,
      anchor=south west,
      label={[anchor=north west, font=\footnotesize\bfseries,
              color=bridgeblue, xshift=8pt, yshift=-6pt]
              north west: Reseau Docker bridge-demo}]
      (net) at (0, 0) {};

% =====================================================================
% Niveau 1 : 4 simulateurs alignes en haut
% =====================================================================
\node[sim] (sim1) at (2.4, 8.6) {sim-monitor\\Philips};
\node[sim] (sim2) at (5.6, 8.6) {sim-glucometer\\Masimo};
\node[sim] (sim3) at (8.8, 8.6) {sim-spirometer\\Welch Allyn};
\node[sim] (sim4) at (12.0, 8.6) {sim-pump\\Drager};

% Etiquettes de protocole sous chaque sim, sans chevaucher
\node[protolbl] at (2.4, 7.85) {TCP HL7v2};
\node[protolbl] at (5.6, 7.85) {MQTT JSON};
\node[protolbl] at (8.8, 7.85) {fichier XML};
\node[protolbl] at (12.0, 7.85) {RS-232};

% =====================================================================
% Niveau 2 : Bus MQTT et volume partage
% =====================================================================
\node[bus] (mqtt) at (5.6, 6.3) {MQTT broker\\eclipse-mosquitto};
\node[vol] (vol)  at (8.8, 6.3) {Volume \texttt{drop/}};

% =====================================================================
% Niveau 3 : Bridge core (centre) et Dashboard (droite)
% =====================================================================
\node[bridgecore] (bridge) at (5.6, 4.0) {IoT Edge Bridge\\Python + FastAPI};
\node[font=\scriptsize\itshape, color=bridgegray, anchor=north]
  at (5.6, 3.4) {Adapter . Parser . Mapper . Transport};

\node[dashbox] (dashb) at (12.0, 4.0) {Dashboard\\:8080};

% =====================================================================
% Niveau 4 : HAPI FHIR Server
% =====================================================================
\node[hapi] (hapi) at (5.6, 1.4) {HAPI FHIR Server\\:8081};

% =====================================================================
% Acteur exterieur (auteur)
% =====================================================================
\node[anchor=west, font=\footnotesize\itshape, color=bridgegray]
  (user) at (16.0, 4.0) {Auteur};

% =====================================================================
% Fleches : simulateurs vers leurs cibles
% =====================================================================
% sim1 (TCP) directement vers Bridge
\draw[flowarrow] (sim1.south) -- ++(0, -0.3)
                 -| ([xshift=-2.0cm]bridge.north);

% sim2 (MQTT) vers MQTT broker
\draw[flowarrow] (sim2.south) -- (mqtt.north);

% sim3 (fichier) vers volume
\draw[flowarrow] (sim3.south) -- (vol.north);

% sim4 (RS-232) directement vers Bridge
\draw[flowarrow] (sim4.south) -- ++(0, -0.3)
                 -| ([xshift=2.0cm]bridge.north);

% MQTT et volume vers Bridge
\draw[flowarrow] (mqtt.south)  -- ([xshift=-0.7cm]bridge.north);
\draw[flowarrow] (vol.south)   -- ([xshift=0.7cm]bridge.north);

% Bridge vers HAPI
\draw[flowarrow] (bridge.south) -- (hapi.north)
  node[midway, right=2pt, font=\scriptsize\itshape, color=bridgegray] {POST FHIR};

% Bridge vers Dashboard
\draw[flowarrow] (bridge.east) -- (dashb.west)
  node[midway, above, font=\scriptsize\itshape, color=bridgegray] {WebSocket};

% =====================================================================
% Acces auteur vers Dashboard (en pointille, hors reseau Docker)
% =====================================================================
\draw[->, dashed, thick, color=bridgegray]
  (user.west) -- (dashb.east)
  node[midway, above, font=\scriptsize\itshape, color=bridgegray] {HTTP};

\end{tikzpicture}
}\\[1mm]
    {\itshape\scriptsize 7 conteneurs orchestr\'es par Docker Compose. \texttt{docker compose up} lance tout.}
  \end{center}
  \note{Sept conteneurs sur un meme reseau Docker. Quatre simulateurs, le bridge, HAPI FHIR, et un broker MQTT. Tout tient sur un poste local. Reproductible a l'identique chez n'importe lequel d'entre vous.}
\end{frame}

% =====================================================================
% F8. STACK TECHNOLOGIQUE
% =====================================================================
\partref{Partie 3/5 : L'architecture}
\chref{Ch.4}
\begin{frame}{Stack technologique}
  \begin{center}
    \footnotesize
    \zebra
    \begin{tabular}{@{}c l l@{}}
      \toprule
      & \textbf{Couche} & \textbf{Technologie} \\
      \midrule
      \textcolor{bridgeblue}{\faPython}        & C\oe ur middleware     & Python 3.12 + FastAPI \\
      \textcolor{bridgeblue}{\faDatabase}      & Persistance locale     & SQLite WAL chiffr\'e \\
      \textcolor{bridgeteal}{\faRandom}        & Dashboard              & HTMX + WebSocket \\
      \textcolor{bridgeteal}{\faRedo}          & Transport s\'ecuris\'e & httpx + Tenacity (retry exponentiel) \\
      \textcolor{bridgeblue}{\faDocker}        & Orchestration          & Docker Compose \\
      \textcolor{bridgegreen}{\faServer}       & Cible FHIR             & HAPI FHIR Server (image officielle) \\
      \bottomrule
    \end{tabular}
  \end{center}
  \vspace{0.4em}
  \centerline{\itshape\scriptsize Que des briques \'eprouv\'ees. Choix volontairement conservateur.}
  \note{Python pour la productivite, FastAPI pour la performance asyncio, SQLite WAL pour la resilience, HTMX pour un dashboard sans framework lourd. Que des briques eprouvees, rien d'exotique. Choix volontairement conservateur.}
\end{frame}

% =====================================================================
% F9. CHAINE CAPTEUR-ACTIONNEUR
% =====================================================================
\partref{Partie 4/5 : Les capteurs}
\chref{Ch.5}
\begin{frame}{La cha\^ine IoT m\'edicale et notre point d'insertion}
  \begin{center}
    \resizebox{0.92\textwidth}{!}{% =====================================================================
% figures/fig-chaine-capteur-actionneur.tex
% Chaine canonique du cours : capteur -> microcontroleur -> communication
%                          -> application -> actionneur
% La passerelle se substitue au maillon Application.
% =====================================================================
\begin{tikzpicture}[
  font=\sffamily,
  every node/.style={font=\footnotesize\sffamily},
  >={Stealth[length=2.5mm]},
  thick,
  block/.style={
    rectangle, rounded corners=2pt,
    draw=bridgeblue, fill=bridgeblue!8, text=bridgeblue,
    minimum width=2.5cm, minimum height=1.1cm, align=center,
    inner sep=3pt, font=\footnotesize\bfseries\sffamily
  },
  highlight/.style={
    rectangle, rounded corners=2pt,
    draw=bridgeteal, fill=bridgeteal!15, text=bridgeteal,
    minimum width=2.5cm, minimum height=1.1cm, align=center,
    inner sep=3pt, font=\footnotesize\bfseries\sffamily
  },
  flow/.style={-{Stealth[length=3mm]}, thick, draw=bridgeblue}
]

% --- Chaine horizontale ----------------------------------------------
\node[block] (capt) at (0, 0) {Capteur};
\node[block] (mcu)  at (3.0, 0) {Micro-\\contr\^oleur};
\node[block] (comm) at (6.0, 0) {Module de\\communication};
\node[highlight] (app) at (9.0, 0) {Application\\(notre passerelle)};
\node[block] (act)  at (12.0, 0) {Actionneur};

\draw[flow] (capt.east) -- (mcu.west);
\draw[flow] (mcu.east)  -- (comm.west);
\draw[flow] (comm.east) -- (app.west);
\draw[flow] (app.east)  -- (act.west);

% --- Annotations sous les blocs --------------------------------------
\node[font=\scriptsize\itshape, color=bridgegray, align=center]
  at ([yshift=-0.85cm]capt.south) {grandeur\\physique};
\node[font=\scriptsize\itshape, color=bridgegray, align=center]
  at ([yshift=-0.85cm]mcu.south) {numerise\\conditionne};
\node[font=\scriptsize\itshape, color=bridgegray, align=center]
  at ([yshift=-0.85cm]comm.south) {RS-232, BLE,\\TCP, fichier};
\node[font=\scriptsize\itshape, color=bridgeteal, align=center]
  at ([yshift=-0.85cm]app.south) {capture passive,\\transformation FHIR};
\node[font=\scriptsize\itshape, color=bridgegray, align=center]
  at ([yshift=-0.85cm]act.south) {pompe,\\alarme, etc.};

% --- Surbrillance "passerelle" ---------------------------------------
\draw[bridgeteal, thick, dashed, rounded corners=4pt]
  ([xshift=-0.3cm, yshift=0.4cm]app.north west) rectangle
  ([xshift=0.3cm, yshift=-0.4cm]app.south east);
\node[anchor=south, font=\scriptsize\bfseries, color=bridgeteal]
  at ([yshift=0.45cm]app.north) {Substitution};

\end{tikzpicture}
}\\[2mm]
    {\itshape\small Notre passerelle se substitue au maillon \emph{application}. Capture passive, jamais d'intrusion mat\'erielle.}
  \end{center}
  \note{Dans la chaine classique capteur, conditionnement, transmission, application, on s'insere uniquement au niveau application. On n'ouvre aucun boitier. C'est cette capture passive qui nous garde compatibles MDR.}
\end{frame}

% =====================================================================
% F10. QUATRE EQUIPEMENTS REPRESENTATIFS
% =====================================================================
\partref{Partie 4/5 : Les capteurs}
\chref{Ch.5}
\begin{frame}{Quatre \'equipements, quatre canaux}
  \begin{center}
    \scriptsize
    \zebra
    \begin{tabular}{@{}c l l l l@{}}
      \toprule
      & \textbf{\'Equipement} & \textbf{Famille} & \textbf{Canal} & \textbf{LOINC} \\
      \midrule
      \textcolor{bridgealert}{\faHeartbeat}    & Philips IntelliVue MX450 & Moniteur multiparam.   & TCP MLLP HL7v2  & SpO2, FC, PA \\
      \textcolor{bridgealert}{\faTint}         & Masimo Radical-7         & Oxym\`etre POC          & BLE \(\to\) MQTT JSON & SpO2, FC \\
      \textcolor{bridgealert}{\faChartLine}    & Welch Allyn CP150        & ECG 12 d\'erivations    & Fichier d\'epos\'e & \'Etude ECG \\
      \textcolor{bridgealert}{\faLungs}        & Dr\"ager Evita V500      & Ventilateur            & RS-232 MEDIBUS  & FiO2, PEEP, VT, FR \\
      \bottomrule
    \end{tabular}\\[2mm]
    {\itshape\small 4 protocoles, 1 seul c\oe ur, 1 profil JSON par dispositif.}
  \end{center}
  \note{Quatre dispositifs, quatre mondes techniques differents. TCP, BLE, fichier, port serie. Le coeur du bridge ne change pas. Seul le profil JSON change. C'est la promesse de la solution.}
\end{frame}

% =====================================================================
% F11. SECURITE ET CONFIDENTIALITE PAR CONSTRUCTION
% =====================================================================
\partref{Partie 5/5 : L'\'ethique}
\chref{Ch.7}
\begin{frame}{S\'ecurit\'e et confidentialit\'e par construction}
  \begin{center}
    \scriptsize
    \zebra
    \begin{tabularx}{\linewidth}{@{}c l l X@{}}
      \toprule
      & \textbf{Pilier} & \textbf{Mise en \oe uvre} & \textbf{Pourquoi c'est l\`a} \\
      \midrule
      \textcolor{bridgeblue}{\faLock}                & Confidentialit\'e & TLS 1.3, minimisation, dashboard sans PII & La donn\'ee patient ne circule jamais en clair ; l'op\'erateur n'a pas besoin de la voir pour faire son m\'etier. \\
      \textcolor{bridgeteal}{\faUserShield}           & Int\'egrit\'e     & SHA-256 sur file de repli, validation FHIR & On d\'etecte corruption ou mapping fautif \emph{avant} d'envoyer une mesure erron\'ee au DPI. \\
      \textcolor{bridgegreen}{\faRedo}               & Disponibilit\'e   & Retry exponentiel, fallback SQLite, hot-reload & Coupure r\'eseau 30 min : aucune perte. Nouveau profil : aucun red\'emarrage. \\
      \textcolor{bridgeblue}{\faClipboardCheck}      & Tra\c{c}abilit\'e & Logs JSON UTC, sans valeur clinique & Audit complet (qui/quand/quoi) sans cr\'eer un nouveau d\'ep\^ot de donn\'ees patient. \\
      \bottomrule
    \end{tabularx}
  \end{center}
  \note{Securite sur les quatre piliers CIA plus tracabilite. Pour chacun, la justification : pourquoi cette mise en oeuvre, et ce qu'elle previent concretement. La minimisation, c'est pas un slogan : le dashboard n'affiche jamais une donnee patient.}
\end{frame}

% =====================================================================
% F12. DEFIS ETHIQUES
% =====================================================================
\partref{Partie 5/5 : L'\'ethique}
\chref{Ch.7}
\begin{frame}{Trois questions \'ethiques, trois r\'eponses}
  \scriptsize
  \zebra
  \begin{tabularx}{\linewidth}{@{}c p{4.0cm} X@{}}
    \toprule
    & \textbf{Question} & \textbf{R\'eponse} \\
    \midrule
    \textcolor{bridgealert}{\faQuestionCircle} &
    \textbf{Responsabilit\'e du mapping.} Qui r\'epond si l'unit\'e est mal cod\'ee~? &
    Profil JSON sign\'e (auteur, version, date). Chaque mesure porte sa r\'ef\'erence. Responsabilit\'e \textbf{nominative}, pas diluee. \\
    \textcolor{bridgealert}{\faQuestionCircle} &
    \textbf{Consentement au partage hub.} Distinct du RGPD g\'en\'eral. &
    Le bridge \emph{produit}, ne \emph{partage pas}. Le consentement BIHR/MetaHub reste g\'er\'e \textbf{par le DPI}. \\
    \textcolor{bridgealert}{\faQuestionCircle} &
    \textbf{Risque de datafication.} Capter facile $\neq$ capter sans limite. &
    Activer un dispositif = d\'ecision explicite, valid\'ee en \textbf{comit\'e clinique}. Pas de d\'eploiement automatique. \\
    \bottomrule
  \end{tabularx}
  \vspace{2mm}
  \begin{bridgewarn}[\faExclamationTriangle~AIPD obligatoire en prod]
    Avant tout d\'eploiement r\'eel : analyse d'impact RGPD formelle, avec consultation du DPO (RGPD art.~35).
  \end{bridgewarn}
  \note{Trois questions qu'on ne peut pas eluder, et trois reponses concretes : signature des profils pour la responsabilite, separation bridge/DPI pour le consentement, comite clinique pour le risque de datafication. La technique ne dispense pas de la reflexion, mais elle peut l'outiller.}
\end{frame}

% =====================================================================
% F13. TRANSITION DEMO
% =====================================================================
\partref{}
\chref{D\'emo}
\begin{frame}[plain]{Place \`a la d\'emo}
  \vspace{1.5cm}
  \begin{center}
    {\Large\bfseries\color{bridgeblue} Maintenant, on lance.}\\[6mm]
    {\itshape\large\ttfamily docker compose up}
  \end{center}
  \note{La theorie c'est bien. Maintenant on lance. Cinq minutes en live, sur Docker. Vous allez voir les quatre simulateurs publier en FHIR sur HAPI en temps reel. Vous allez voir le dashboard passer au vert.}
\end{frame}

% =====================================================================
% F14. CONCLUSION 1 : "L'HISTOIRE QUE JE VOUS AI RACONTEE"
% =====================================================================
\partref{}
\chref{}
\begin{frame}{L'histoire que je vous ai racont\'ee}

\begin{itemize}[leftmargin=*, itemsep=0.4em]
  \item \textcolor{bridgealert}{\faExclamationTriangle}~\textbf{Au d\'ebut} : dispositifs muets, donn\'ees captives.
  \item \textcolor{bridgeblue}{\faProjectDiagram}~\textbf{Nous avons pos\'e} un middleware configurable, transparent, MDR/RGPD-aware.
  \item \textcolor{bridgeteal}{\faRoute}~\textbf{Quatre couches} : Adapter, Parser, Mapper, Transport. Un seul c\oe ur.
  \item \textcolor{bridgeblue}{\faUserMd}~\textbf{Une supervision OT} biom\'edicale, sans jamais exposer la donn\'ee patient.
  \item \textcolor{bridgegreen}{\faNetworkWired}~\textbf{Une ouverture aux plugins}, pour qu'aucun \'editeur seul ne soit le goulet.
\end{itemize}

\vspace{0.6em}
\centerline{\itshape\large\color{bridgeblue} Le chaos est devenu de l'ordre. Sans toucher au mat\'eriel.}
\note{Slide 14. C'est le moment de clore le recit. Premiere conclusion : narrative. On rappelle l'arc complet.}
\end{frame}

% =====================================================================
% F15. CONCLUSION 2 : "L'HISTOIRE EN UNE SLIDE"
% =====================================================================
\partref{}
\chref{}
\begin{frame}{L'histoire en une slide}

\begin{itemize}[label={}, leftmargin=*, itemsep=0.5em]
  \item \textcolor{bridgealert}{\faHospital}~\textbf{1.~Le contexte.} Mesures captives, BIHR attendu, MDR verrou.
  \item \textcolor{bridgeteal}{\faProjectDiagram}~\textbf{2.~Le syst\`eme.} 4 sources, 1 c\oe ur, FHIR R4 en sortie.
  \item \textcolor{bridgeblue}{\faPython}~\textbf{3.~L'architecture.} Adapter, Parser, Mapper, Transport (Docker, Python).
  \item \textcolor{bridgealert}{\faHeartbeat}~\textbf{4.~Les capteurs.} 4 protocoles, 4 \'equipements, 8 codes LOINC.
  \item \textcolor{bridgeblue}{\faUserShield}~\textbf{5.~L'\'ethique.} CIA + minimisation + AIPD obligatoire en prod.
\end{itemize}

\note{Slide 15. Seconde conclusion : synthese pedagogique, pour ancrer dans la memoire les 5 piliers du pitch.}
\end{frame}

% =====================================================================
% F16. OUVERTURE : OU EN SOMMES-NOUS
% =====================================================================
\partref{Ouverture sur le futur}
\chref{Ch.6}
\begin{frame}{Et apr\`es ? L\`a o\`u nous en sommes}

\scriptsize
\zebra
\begin{tabularx}{\linewidth}{@{}c p{3.4cm} X@{}}
\toprule
& \textbf{Statut} & \textbf{Pourquoi} \\
\midrule
\textcolor{bridgeblue}{\faBalanceScale} & \textbf{Non-DM} (middleware IT) &
Le bridge \textbf{transporte} et \textbf{traduit}, il ne mesure ni ne diagnostique. MDR 2017/745 art.~2 ne le qualifie donc pas de DM. \\
\textcolor{bridgegreen}{\faCheckCircle} & \textbf{Niveau M2} mhealth.belgium &
M1 info, \textbf{M2 parcours coordonn\'e}, M3 rembours\'e. M2 = palier r\'ealiste pour un middleware non rembours\'e INAMI. \\
\bottomrule
\end{tabularx}

\vspace{0.4em}
\scriptsize
\zebra
\begin{tabularx}{\linewidth}{@{}c X r@{}}
\toprule
& \textbf{Indicateur} & \textbf{Cible} \\
\midrule
\textcolor{bridgeteal}{\faClock}      & Latence p95 ingestion (capteur $\to$ HAPI) & $< 500$ ms \\
\textcolor{bridgegreen}{\faCheckCircle} & Taux de transmission r\'eussie & $> 99~\%$ \\
\textcolor{bridgeblue}{\faUserMd}     & Min/soignant/jour \'economis\'ees & $> 30$ \\
\textcolor{bridgeblue}{\faPython}     & LdC pour ajouter un dispositif & 0 \\
\bottomrule
\end{tabularx}

\note{Slide 16. Statut non-DM justifie par l'absence de mesure et de diagnostic. M2 justifie comme objectif realiste pour un middleware hospitalier non rembourse. Indicateurs concrets en bas.}
\end{frame}

% =====================================================================
% F17. OUVERTURE : ROADMAP F1/F2/F3
% =====================================================================
\partref{Ouverture sur le futur}
\chref{Roadmap}
\begin{frame}{Et apr\`es ? Trois \'evolutions \`a venir}

\begin{itemize}[leftmargin=*, itemsep=0.5em]
  \item \textcolor{bridgegreen}{\faPlug}~\textbf{F1.} Plugins communautaires (mod\`ele HACS, npm).
  \item \textcolor{bridgeteal}{\faChartLine}~\textbf{F2.} D\'etection d'anomalies en bord (IA l\'eg\`ere ONNX).
  \item \textcolor{bridgeblue}{\faNetworkWired}~\textbf{F3.} F\'ed\'eration europ\'eenne EHDS (mobilit\'e transfrontali\`ere).
\end{itemize}

\vspace{0.6em}
\begin{bridgeinfo}[\faNetworkWired~Effet de catalogue]
Aucun \'editeur seul ne couvre tout le parc mondial. La communaut\'e \'ecrit ses plugins. Le bridge devient un \textbf{standard de fait}.
\end{bridgeinfo}
\note{Slide 17. Deuxieme prolongement : ce qu'on imagine pour la suite, et la dynamique communautaire qui rend cette suite plausible.}
\end{frame}

% =====================================================================
% F18. MERCI / Q&R
% =====================================================================
\partref{}
\chref{}
\begin{frame}[plain]{Merci}
  \vspace{1.5cm}
  \begin{center}
    {\Large\color{bridgeblue} IoT Edge Bridge}\\[4mm]
    {\large\itshape Questions, remarques, id\'ees de plugins ?}
  \end{center}
  \note{Merci pour votre attention. Si vous avez des questions, des idees d'equipements a integrer, des plugins a ecrire, je suis preneur. Et si vous voulez tester chez vous, le repo est ouvert.}
\end{frame}

\end{document}
